\input{../tex-inputs/latex-leseansicht-vorspann}

\section[Arthur Schnitzler an Gustav Schwarzkopf, 9. 9. 1919]{L04527 Arthur Schnitzler an Gustav Schwarzkopf, 9. 9. 1919}
\nopagebreak\mylabel{L04527v}
\rehead{ }\normalsize\beginnumbering\briefempfaengerindex{Schwarzkopf, Gustav@\textsc{Schwarzkopf, Gustav}!zzzSchnitzler, Arthur@\emph{von Arthur Schnitzler}!1919-09-091@{9. 9. 1919}|(be}
\toendnotes[C]{\smallbreak\pagebreak[2]}
\correspDesc{Versand  durch Arthur Schnitzler am 9. 9. 1919 in Reichenau an der Rax
\newline{}Übermittlung  am 10. 9. 1919 in Reichenau an der Rax
\newline{}Erhalt  durch Gustav Schwarzkopf im Zeitraum [10. 9. 1919 – 14. 9. 1919?] in Wien}\toendnotes[C]{\smallbreak}
\Standort{DLA, A:Schnitzler, HS.1985.1.1897.}
\physDesc{Bildpostkarte, 82 Zeichen
\newline{}Handschrift: Bleistift, deutsche Kurrent
\newline{}Versand: Stempel: »\nobreak{}\oindex{Reichenau an der Rax@\textbf{Reichenau an der Rax}, \emph{Verwaltungsgebiet}|pwk}Reiche\textcolor{gray}{nau}, 10. \textcolor{gray}{IX. 19}, \nobreak{}«.  }\pstart{}{\pb}\textsc{Hrn Gustav Schwarzkopf}\pend{}\pstart{}\textsc{Wien I}\oindex{I., Innere Stadt@\textbf{I., Innere Stadt}, \emph{Verwaltungsgebiet}|pw}\pend{}\pstart{}\textsc{Tiefer Graben 17}\oindex{Wien@\textbf{Wien}!I., Innere Stadt@\textbf{I., Innere Stadt}!Tiefer Graben 17@\textbf{Tiefer Graben 17}, \emph{Wohngebäude}|pw}\pend{}{\bigskip}
\pstart
           \noindent{}\centering{}{\pb}\textcolor{gray}{\textbf{Reichenau (N.-Ö.)\oindex{Reichenau an der Rax@\textbf{Reichenau an der Rax}, \emph{Verwaltungsgebiet}|pw} mit Raxalpe\oindex{Raxalpe@\textbf{Raxalpe}, \emph{Alm}|pw}.}}\pend
           \vspace{1em}
\pstart
           \raggedleft{}{\pb}9. 9. 1919\pend
           \vspace{0.5em}
\pstart
           Herzliche Grüße –!\pend
           
\pstart
           Ihr{\\[\baselineskip]}\spacefill\mbox{Arthur.}\pend
           \leftskip=0em{}\selectlanguage{ngerman}\endnumbering\briefempfaengerindex{Schwarzkopf, Gustav@\textsc{Schwarzkopf, Gustav}!zzzSchnitzler, Arthur@\emph{von Arthur Schnitzler}!1919-09-091@{9. 9. 1919}|)be}\mylabel{L04527h}
\begin{anhang}
\end{anhang}\newcommand{\dateiname}{L04527}\newcommand{\titel}{Arthur Schnitzler an Gustav Schwarzkopf, 9. 9. 1919}\newcommand{\editorInnen}{Herausgegeben von Jahnke, SelmaMüller, Martin Anton}\input{../tex-inputs/latex-leseansicht-abspann}
